%!TEX program = xelatex
%!TEX root = ../all_beamer.tex

%% Ensemble Methods – Beamer Presentation
%% Based on notes by Giorgio Gambosi, University of Rome "Tor Vergata"
%% Style modelled on ml-dectrees-beamer.tex

%% !TEX root = all_beamer.tex
% ==============================================================================
% HEADER LATEX PER PRESENTAZIONI BEAMER
% Corso di Machine Learning - Giorgio Gambosi
% Università di Roma "Tor Vergata"
% ==============================================================================

% ------------------------------------------------------------------------------
% CLASSE DOCUMENTO E METADATI
% ------------------------------------------------------------------------------
\documentclass[aspectratio=169, 10pt]{beamer}

\subtitle{Course of Machine Learning}
\author{Giorgio Gambosi}
\institute[Tor Vergata]{Master's Degree in Computer Science\\University of Rome ``Tor Vergata''}
\date{A.Y. 2025--2026}

% ------------------------------------------------------------------------------
% PACCHETTI MATEMATICI
% ------------------------------------------------------------------------------
\usepackage{amsmath, amssymb, amsfonts, mathtools, amsthm}
\usepackage{bm}              % Simboli matematici in grassetto
\usepackage{upgreek}         % Lettere greche dritte
\usepackage{derivative}      % Gestione semplificata delle derivate

% ------------------------------------------------------------------------------
% PACCHETTI GRAFICI E VISUALIZZAZIONE
% ------------------------------------------------------------------------------
\usepackage{graphicx}        % Inclusione immagini
\usepackage{xcolor}          % Gestione colori
\usepackage{xkcdcolors}      % Colori aggiuntivi stile xkcd
\usepackage{microtype}       % Miglioramenti tipografici

% ------------------------------------------------------------------------------
% PACCHETTI PER TABELLE
% ------------------------------------------------------------------------------
\usepackage{booktabs}        % Tabelle professionali
\usepackage{array}           % Estensioni per tabelle
\usepackage{multirow}        % Celle multi-riga
\usepackage{multicol}        % Layout multi-colonna

% ------------------------------------------------------------------------------
% TIKZ E PGFPLOTS
% ------------------------------------------------------------------------------
\usepackage{tikz}
\usepackage{tikz-network}
\usepackage[customcolors]{hf-tikz}
\usepackage{pgfplots}
\pgfplotsset{compat=1.18}
\usepgfplotslibrary{fillbetween}

% Librerie TikZ
\usetikzlibrary{arrows, arrows.meta, decorations.pathmorphing, decorations.pathreplacing}
\usetikzlibrary{decorations.markings, snakes, positioning, backgrounds, fit, shadows}
\usetikzlibrary{automata, patterns, calc, intersections, through, shapes}
\usetikzlibrary{shapes.geometric, shapes.arrows, scopes, chains, matrix}

% ------------------------------------------------------------------------------
% ALTRI PACCHETTI
% ------------------------------------------------------------------------------
\usepackage[skins]{tcolorbox}  % Box colorati
\usepackage{subcaption}         % Sottocaption per figure
\usepackage{listings}           % Inclusione codice sorgente

% ------------------------------------------------------------------------------
% TEMA E STILE BEAMER
% ------------------------------------------------------------------------------
\usetheme{Madrid}
\usecolortheme{whale}
\useinnertheme{rounded}
\usefonttheme{professionalfonts}
\setbeamertemplate{navigation symbols}{}  % Rimuove simboli di navigazione

% ------------------------------------------------------------------------------
% DEFINIZIONE COLORI PERSONALIZZATI
% ------------------------------------------------------------------------------
% Colori principali del tema
\definecolor{MainBlue}{RGB}{31,56,100}      % Blu principale (header, titoli)
\definecolor{AccentBlue}{RGB}{70,130,180}   % Blu accento (blocchi)
\definecolor{AlertRed}{RGB}{160,30,30}      % Rosso per alert
\definecolor{ExGreen}{RGB}{30,120,60}       % Verde per esempi
\definecolor{BlockBg}{RGB}{220,232,246}     % Sfondo blocchi
\definecolor{torgray}{RGB}{220,225,235}     % Grigio Tor Vergata

% Colori alternativi (Tor Vergata style)
\definecolor{TvDark}{RGB}{30,50,80}         % Blu scuro header
\definecolor{TvMid}{RGB}{70,100,140}        % Blu medio blocchi
\definecolor{TvLight}{RGB}{210,220,235}     % Blu chiaro sfondi
\definecolor{TvAlert}{RGB}{160,30,30}       % Rosso alert
\definecolor{TvGreen}{RGB}{20,100,50}       % Verde esempi

% Colori supplementari
\definecolor{torred}{RGB}{153,0,0}
\definecolor{torcyan}{RGB}{0,112,192}
\definecolor{cerise}{HTML}{DA3A62}
\definecolor{teal}{HTML}{029386}
\definecolor{slateblue}{HTML}{5B5EA6}
\definecolor{orange}{HTML}{F97306}

% Colori per nodi TikZ
\definecolor{nodeblue}{RGB}{101,140,196}    % xkcd Faded Blue

% ------------------------------------------------------------------------------
% CONFIGURAZIONE COLORI BEAMER
% ------------------------------------------------------------------------------
\setbeamercolor{frametitle}{bg=MainBlue, fg=white}
\setbeamercolor{title}{bg=MainBlue, fg=white}
\setbeamercolor{structure}{fg=MainBlue}

% Blocchi standard
\setbeamercolor{block title}{bg=AccentBlue, fg=white}
\setbeamercolor{block body}{bg=BlockBg}

% Blocchi alert
\setbeamercolor{block title alerted}{bg=AlertRed, fg=white}
\setbeamercolor{block body alerted}{bg=red!8}

% Blocchi esempio
\setbeamercolor{block title example}{bg=ExGreen, fg=white}
\setbeamercolor{block body example}{bg=green!7}

% ------------------------------------------------------------------------------
% CONFIGURAZIONE FONT BEAMER
% ------------------------------------------------------------------------------
\setbeamerfont{frametitle}{size=\normalsize,series=\bfseries}
\setbeamerfont{block title}{size=\small,series=\bfseries}

% ------------------------------------------------------------------------------
% STILI TIKZ PERSONALIZZATI
% ------------------------------------------------------------------------------
\tikzset{
  >=latex,  % Frecce stile LaTeX di default
  % Nodi per reti neurali
  node/.style={thick,circle,draw=myblue,minimum size=22,inner sep=0.5,outer sep=0.6},
  node in/.style={node,green!20!black,draw=mygreen!30!black,fill=mygreen!25},
  node hidden/.style={node,blue!20!black,draw=myblue!30!black,fill=myblue!20},
  node convol/.style={node,orange!20!black,draw=myorange!30!black,fill=myorange!20},
  node out/.style={node,red!20!black,draw=myred!30!black,fill=myred!20},
  % Connessioni
  connect/.style={thick,mydarkblue},
  connect arrow/.style={-{Latex[length=4,width=3.5]},thick,mydarkblue,shorten <=0.5,shorten >=1},
  % Stili numerati per facile mappatura
  node 1/.style={node in},
  node 2/.style={node hidden},
  node 3/.style={node out}
}

% ------------------------------------------------------------------------------
% AMBIENTI TEOREMI
% ------------------------------------------------------------------------------
\setbeamertemplate{theorems}[numbered]
\newtheorem{mytheorem}{Theorem}
\newtheorem{mydefinition}{Definition}

% ------------------------------------------------------------------------------
% CONFIGURAZIONE LISTINGS (CODICE)
% ------------------------------------------------------------------------------
\lstset{
  language=Python,
  basicstyle=\ttfamily\tiny,
  keywordstyle=\color{MainBlue}\bfseries,
  commentstyle=\color{gray},
  frame=single,
  breaklines=true,
}

% ==============================================================================
% MACRO MATEMATICHE PERSONALIZZATE
% ==============================================================================

% ------------------------------------------------------------------------------
% SPAZI CALLIGRAFICI (Calligraphic spaces)
% ------------------------------------------------------------------------------
\providecommand{\calX}{\mathcal{X}}         % Spazio input
\providecommand{\calY}{\mathcal{Y}}         % Spazio output
\providecommand{\calH}{\mathcal{H}}         % Spazio ipotesi
\providecommand{\calT}{\mathcal{T}}         % Training set
\providecommand{\calA}{\mathcal{A}}         % Algoritmo
\providecommand{\calR}{\mathcal{R}}         % Rischio
\providecommand{\calF}{\mathcal{F}}         % Famiglia di funzioni
\providecommand{\calL}{\mathcal{L}}         % Loss function

% Alias per compatibilità
\providecommand{\Rr}{\mathcal{R}}
\providecommand{\Tr}{\mathcal{T}}
\providecommand{\Hr}{\mathcal{H}}

% ------------------------------------------------------------------------------
% RISCHIO E PROBABILITÀ
% ------------------------------------------------------------------------------
\providecommand{\barR}{\overline{\mathcal{R}}}  % Rischio empirico
\providecommand{\EmpR}{\overline{\mathcal{R}}}  % Rischio empirico (alias)
\providecommand{\Rbar}{\overline{\mathcal{R}}}  % Rischio empirico (alias)
\providecommand{\pM}{p_M}                       % Probabilità modello
\providecommand{\pC}{p_C}                       % Probabilità corretta
\providecommand{\Prob}[2]{\mathbb{P}_{#1}\!\left[#2\right]}  % Probabilità

% ------------------------------------------------------------------------------
% FUNZIONI INDICATRICI
% ------------------------------------------------------------------------------
\providecommand{\indicator}[1]{\mathbf{1}\!\left[#1\right]}  % Funzione indicatrice
\providecommand{\ind}[1]{\mathbf{1}[#1]}                      % Versione breve

% ------------------------------------------------------------------------------
% OPERATORI MATEMATICI
% ------------------------------------------------------------------------------
\providecommand{\argmax}{\operatorname*{arg\,max}}  % Argomento del massimo
\providecommand{\argmin}{\operatorname*{arg\,min}}  % Argomento del minimo
\providecommand{\vcd}[1]{\operatorname{VCdim}(#1)} % VC dimension
\providecommand{\defeq}{\stackrel{\text{def}}{=}}  % Definito come

% ------------------------------------------------------------------------------
% VETTORI E MATRICI (Bold Math)
% ------------------------------------------------------------------------------
% Vettori minuscoli
\providecommand{\bx}{\mathbf{x}}            % Vettore x
\providecommand{\bz}{\mathbf{z}}            % Vettore z
\providecommand{\bw}{\mathbf{w}}            % Vettore peso
\providecommand{\btt}{\mathbf{t}}           % Vettore target
\providecommand{\wvec}{\mathbf{w}}          % Vettore peso (alias)
\providecommand{\bth}{\boldsymbol{\theta}}  % Parametro theta

% Matrici maiuscole
\providecommand{\bX}{\mathbf{X}}            % Matrice dati
\providecommand{\bZ}{\mathbf{Z}}            % Matrice Z
\providecommand{\Xmat}{\mathbf{X}}          % Matrice dati (alias)
\providecommand{\Hmat}{\mathbf{H}}          % Matrice H
\providecommand{\Smat}{\mathbf{S}}          % Matrice scatter generica
\providecommand{\bS}{\mathbf{S}}            % Matrice scatter
\providecommand{\Ii}{\mathbf{I}}            % Matrice identità
\providecommand{\bZero}{\mathbf{0}}         % Vettore zero
\providecommand{\bPsi}{\boldsymbol{\Psi}}   % Matrice Psi

% Medie (overline)
\providecommand{\ox}{\overline{\mathbf{x}}} % Media vettore x
\providecommand{\ow}{\overline{\mathbf{w}}} % Media vettore w
\providecommand{\oX}{\overline{\mathbf{X}}} % Media matrice X

% Matrici scatter specifiche (analisi discriminante)
\providecommand{\SW}{\mathbf{S}_W}          % Within-class scatter
\providecommand{\SB}{\mathbf{S}_B}          % Between-class scatter
\providecommand{\ST}{\mathbf{S}_T}          % Total scatter
\providecommand{\GX}{G_{\mathbf{X}}}        % Matrice Gram

% ------------------------------------------------------------------------------
% COMANDI DI FORMATTAZIONE TESTO
% ------------------------------------------------------------------------------
\providecommand{\term}[1]{\textbf{\color{accent}#1}}        % Termine importante
\providecommand{\halert}[1]{{\color{AlertRed}\textbf{#1}}}  % Alert rosso
\providecommand{\mlalert}[1]{{\color{AccentBlue}\textbf{#1}}}  % Alert blu
\providecommand{\mlemph}[1]{{\color{MainBlue}\textit{#1}}}  % Enfasi blu

% ==============================================================================
% FINE HEADER
% ==============================================================================

\newcommand{\blankpage}{
\clearpage
\thispagestyle{empty}
\null
\clearpage
}

%
%%% ── Packages ────────────────────────────────────────────────────────────────
%\usepackage{amsmath,amssymb,bm}
%\usepackage{booktabs}
%\usepackage{graphicx}
%\usepackage{xcolor}
%\usepackage{xkcdcolors}
%\usepackage{tikz}
%\usepackage{subcaption}
%
%%% ── Math shorthands ─────────────────────────────────────────────────────────
%\newcommand{\SW}{\mathbf{S}_W}
%\newcommand{\SB}{\mathbf{S}_B}
%\newcommand{\ST}{\mathbf{S}_T}
%\newcommand{\Smat}{\mathbf{S}}
%\newcommand{\GX}{G_{\mathbf{X}}}

%% ── Metadata ────────────────────────────────────────────────────────────────
\title[Ensemble Methods]{Ensemble Methods}
%\subtitle{Course of Machine Learning}
%\author{Giorgio Gambosi}
%\institute[Tor Vergata]{Master's Degree in Computer Science\\University of Rome ``Tor Vergata''}
%\date{A.Y. 2025--2026}

%% ════════════════════════════════════════════════════════════════════════════
%\begin{document}

\begin{frame}
  \titlepage
\end{frame}

\begin{frame}{Outline}
  \tableofcontents
\end{frame}

%% ════════════════════════════════════════════════════════════════════════════
\section{Introduction to Ensemble Methods}

\begin{frame}{What Are Ensemble Methods?}
  \framesubtitle{Combining multiple learners for superior prediction}
  \begin{itemize}
    \item An \alert{ensemble method} builds several predictive models and combines their outputs
          to obtain a final predictor that is more accurate, more stable, or both.
    \item Core motivation: a single model is limited by the particular sample on which it was
          trained and by the inductive biases of its hypothesis class.
    \item By \alert{aggregating multiple models}, we reduce the impact of these limitations,
          much like averaging repeated measurements reduces noise.
    \item From a probabilistic viewpoint: ensembles reduce \alert{uncertainty} by combining
          multiple plausible rules, each affected by sampling variability.
  \end{itemize}
  \begin{block}{Error correction through diversity}
    Idiosyncratic mistakes made by one model may be compensated by other models, while
    consistent patterns across models are \alert{reinforced}.
  \end{block}
\end{frame}

\begin{frame}{Two Fundamental Paradigms}
  \begin{columns}[t]
    \column{0.48\textwidth}
      \begin{block}{Bagging (parallel)}
        \begin{itemize}
          \item Constructs a \alert{committee} of $L$ models trained independently on
                perturbed versions of the data.
          \item Aggregates predictions through voting or averaging.
          \item Primarily reduces \alert{variance}.
          \item Models are diverse through data perturbation.
        \end{itemize}
      \end{block}
    \column{0.48\textwidth}
      \begin{block}{Boosting (sequential)}
        \begin{itemize}
          \item Constructs models \alert{sequentially}, each trained to address the
                shortcomings of the current ensemble.
          \item Later models focus on hard examples.
          \item Primarily reduces \alert{bias}.
          \item Diversity through adaptive reweighting.
        \end{itemize}
      \end{block}
  \end{columns}
  \vspace{0.6em}
  \begin{itemize}
    \item Both produce ensembles, but differ in how diversity is created and what aspect
          of error they primarily target.
  \end{itemize}
\end{frame}

%% ════════════════════════════════════════════════════════════════════════════
\section{Bootstrap}

\begin{frame}{The Bootstrap Principle}
  \framesubtitle{Resampling to approximate sampling uncertainty}
  \begin{itemize}
    \item The \alert{bootstrap} is a fundamental resampling tool whose purpose is to approximate
          sampling uncertainty when the true data-generating distribution is unknown.
    \item Key observation: in most realistic problems, the underlying distribution $\mathcal{F}$
          from which data are drawn cannot be accessed directly.
    \item Bootstrap solution: treat the observed dataset itself as a \alert{proxy for $\mathcal{F}$},
          and resample from it repeatedly to mimic collecting new data.
  \end{itemize}
  \begin{block}{Converting a hard problem into a computational one}
    The bootstrap converts reasoning about the sampling distribution under an unknown distribution
    into a \alert{computational problem} that can be addressed by simulation.
  \end{block}
\end{frame}

\begin{frame}{The Empirical Distribution}
  \framesubtitle{The plug-in principle}
  \begin{itemize}
    \item Given training data $(\mathbf{x}_i, t_i)$ for $i = 1, \ldots, n$ drawn i.i.d.\ from
          unknown joint distribution $\mathcal{F}$.
    \item Bootstrap replaces $\mathcal{F}$ with the \alert{empirical distribution}
          $\hat{\mathcal{F}}$, which assigns equal probability mass $1/n$ to each observed
          training pair.
  \end{itemize}
  \begin{block}{Empirical probability mass function}
    \[
    \hat{p}(\mathbf{x}, t) =
    \begin{cases}
    \dfrac{1}{n} & \text{if } \exists\, i : (\mathbf{x}, t) = (\mathbf{x}_i, t_i), \\
    0            & \text{otherwise.}
    \end{cases}
    \]
  \end{block}
  \begin{itemize}
    \item This is a discrete distribution supported entirely on the observed sample.
    \item As $n$ grows, $\hat{\mathcal{F}}$ converges to $\mathcal{F}$ by the
          \alert{Glivenko--Cantelli theorem}.
  \end{itemize}
\end{frame}

\begin{frame}{Bootstrap Sampling Procedure}
  \begin{itemize}
    \item To draw a \alert{bootstrap sample} of size $n$, we sample $n$ times \emph{with replacement}
          from the original dataset.
    \item Each draw selects one of the $n$ observed pairs with equal probability $1/n$,
          independently of all other draws.
    \item Because sampling is \alert{with replacement}, some original observations appear
          multiple times in a bootstrap sample, while others do not appear at all.
  \end{itemize}
  \begin{block}{Bootstrap replication}
    Repeating this procedure $B$ times produces $B$ distinct bootstrap datasets
    $\mathcal{D}^{(1)}, \mathcal{D}^{(2)}, \ldots, \mathcal{D}^{(B)}$, each of size $n$ but
    with different compositions of training examples.
  \end{block}
  \begin{itemize}
    \item Each bootstrap dataset is used to train a separate model, yielding $B$ fitted models.
    \item The distribution of these models approximates the distribution of models we would
          obtain by repeatedly collecting new samples from $\mathcal{F}$.
  \end{itemize}
\end{frame}

\begin{frame}{Expected Coverage in a Bootstrap Sample}
  \begin{itemize}
    \item Consider a particular training example $(\mathbf{x}_i, t_i)$ from the original dataset.
    \item In one bootstrap sample, the probability that this example is \alert{not selected}
          in a single draw is $(1 - 1/n)$.
    \item Since we make $n$ independent draws, the probability that the example is never
          selected in the entire bootstrap sample is:
  \end{itemize}
  \begin{block}{Probability of exclusion}
    \[
    \left(1 - \frac{1}{n}\right)^n \xrightarrow{n \to \infty} e^{-1} \approx 0.368.
    \]
  \end{block}
  \begin{itemize}
    \item Therefore, approximately \alert{63.2\%} of the original training examples appear
          at least once in any given bootstrap sample.
    \item The remaining \alert{36.8\%} are \emph{out-of-bag} (OOB) examples for that sample.
  \end{itemize}
\end{frame}

\begin{frame}{Out-of-Bag Examples and Validation}
  \begin{itemize}
    \item For each bootstrap sample $\mathcal{D}^{(b)}$, the examples not selected constitute
          the \alert{out-of-bag (OOB)} set for that bootstrap replicate.
    \item These OOB examples were not used to train the model, so they provide an
          \alert{unbiased validation set} for evaluating that model's performance.
  \end{itemize}
  \begin{block}{OOB error estimation}
    By aggregating OOB predictions across all bootstrap replicates, we can estimate
    generalization error \alert{without setting aside a separate validation set}.
  \end{block}
  \begin{itemize}
    \item This is particularly valuable when data is scarce: the full dataset is used for
          training, yet we still obtain honest performance estimates.
    \item OOB evaluation is a distinctive feature of bootstrap-based ensembles like bagging
          and random forests.
  \end{itemize}
\end{frame}

%% ════════════════════════════════════════════════════════════════════════════
\section{Bagging}

\begin{frame}{Bootstrap Aggregation (Bagging)}
  \framesubtitle{Variance reduction through model averaging}
  \begin{itemize}
    \item \alert{Bagging} (Bootstrap AGGregatING) applies the bootstrap principle directly
          to supervised learning.
    \item Given a training set $\mathcal{D}$ of size $n$, we generate $B$ bootstrap samples
          $\mathcal{D}^{(1)}, \ldots, \mathcal{D}^{(B)}$.
    \item On each bootstrap sample, we train a \alert{base learner} $\hat{f}^{(b)}$, producing
          $B$ fitted models.
    \item Final prediction is obtained by \alert{aggregating} the predictions of all $B$ models.
  \end{itemize}
  \begin{block}{Aggregation rules}
    \begin{itemize}
      \item \textbf{Regression}: average the predictions:
            $\displaystyle \hat{f}_{\text{bag}}(\mathbf{x}) = \frac{1}{B}\sum_{b=1}^{B} \hat{f}^{(b)}(\mathbf{x})$.
      \item \textbf{Classification}: majority vote among the predicted class labels.
    \end{itemize}
  \end{block}
\end{frame}

\begin{frame}{Why Does Bagging Reduce Variance?}
  \begin{itemize}
    \item Consider the case of regression with squared loss. The expected prediction error of
          a single model $\hat{f}$ can be decomposed as:
          \[
          \mathbb{E}[(t - \hat{f}(\mathbf{x}))^2] = \text{Bias}^2 + \text{Variance} + \text{Irreducible error}.
          \]
    \item If the base learner has high variance (e.g., deep decision trees), individual predictions
          are highly sensitive to which particular sample was used for training.
    \item Averaging reduces variance: if predictions $\hat{f}^{(1)}(\mathbf{x}), \ldots, \hat{f}^{(B)}(\mathbf{x})$
          were independent with variance $\sigma^2$, the variance of the average would be $\sigma^2/B$.
  \end{itemize}
  \begin{alertblock}{In practice}
    Bootstrap samples are \alert{not independent} (they overlap), so variance reduction is less
    than $1/B$. Nevertheless, substantial reduction is typically achieved.
  \end{alertblock}
\end{frame}

\begin{frame}{Bagging Does Not Reduce Bias}
  \begin{itemize}
    \item Bagging primarily addresses \alert{variance}, not bias.
    \item Each bootstrap sample is still drawn from the empirical distribution $\hat{\mathcal{F}}$,
          which is an approximation of the true distribution $\mathcal{F}$.
    \item The bias of the bagged predictor is approximately the same as the bias of a single
          base learner trained on the full dataset.
  \end{itemize}
  \begin{block}{Implication}
    Bagging is most effective when combined with \alert{low-bias, high-variance} base learners,
    such as:
    \begin{itemize}
      \item Deep, unpruned decision trees.
      \item Complex models that can overfit easily.
    \end{itemize}
  \end{block}
  \begin{itemize}
    \item If the base learner is already high-bias (e.g., shallow tree), bagging provides
          little benefit because there is little variance to reduce.
  \end{itemize}
\end{frame}

\begin{frame}{Base Learners for Bagging}
  \begin{itemize}
    \item In principle, \alert{any supervised learning algorithm} can serve as the base learner
          for bagging.
    \item However, bagging is most beneficial when the base learner is:
      \begin{itemize}
        \item \alert{Unstable}: small changes in training data lead to large changes in the
              learned model.
        \item \alert{Low bias}: capable of fitting complex patterns when given enough data.
      \end{itemize}
    \item Decision trees are the canonical choice because they satisfy both properties.
  \end{itemize}
  \begin{block}{Stable vs.\ unstable learners}
    \begin{itemize}
      \item Stable learners (e.g., linear regression, nearest neighbors with large $k$):
            bagging provides little improvement.
      \item Unstable learners (e.g., decision trees, neural networks): bagging can yield
            substantial gains.
    \end{itemize}
  \end{block}
\end{frame}

\begin{frame}{Out-of-Bag Error in Bagging}
  \begin{itemize}
    \item For each training example $(\mathbf{x}_i, t_i)$, approximately $36.8\%$ of the
          bootstrap samples will not include it.
    \item These are the \alert{out-of-bag (OOB)} samples for example $i$.
    \item We can predict $t_i$ using only those models for which $(\mathbf{x}_i, t_i)$
          was out-of-bag.
    \item By aggregating OOB predictions across all training examples, we obtain an unbiased
          estimate of generalization error.
  \end{itemize}
  \begin{block}{Advantage}
    OOB error provides a \alert{free cross-validation estimate} without the need for a
    separate validation set or multiple training runs.
  \end{block}
  \begin{itemize}
    \item This is particularly valuable in bagging and random forests, where it has become
          a standard diagnostic.
  \end{itemize}
\end{frame}

\begin{frame}{Number of Bootstrap Replicates}
  \begin{itemize}
    \item The parameter $B$ (number of bootstrap samples) controls the size of the ensemble.
    \item Unlike many hyperparameters, increasing $B$ does \alert{not cause overfitting}.
    \item The bagged predictor converges as $B \to \infty$ to the idealized expectation over
          all possible bootstrap samples from $\hat{\mathcal{F}}$.
    \item In practice, $B$ is chosen based on computational budget rather than statistical
          considerations.
  \end{itemize}
  \begin{block}{Practical recommendation}
    \begin{itemize}
      \item Start with $B = 100$ to $500$ for exploratory analysis.
      \item For production models, $B = 1000$ or more is common.
      \item Monitor OOB error as $B$ increases; it typically stabilizes after a few hundred trees.
    \end{itemize}
  \end{block}
\end{frame}

\begin{frame}{Bagging: Summary of Key Properties}
  \begin{columns}[t]
    \column{0.48\textwidth}
      \begin{block}{Strengths}
        \begin{itemize}
          \item \alert{Reduces variance} substantially for unstable learners.
          \item Simple and easy to parallelize.
          \item Provides OOB error estimates.
          \item Does not overfit as $B$ increases.
        \end{itemize}
      \end{block}
    \column{0.48\textwidth}
      \begin{block}{Limitations}
        \begin{itemize}
          \item Does \alert{not reduce bias}.
          \item Less effective with stable learners.
          \item Loss of interpretability compared to single model.
          \item Requires computational resources proportional to $B$.
        \end{itemize}
      \end{block}
  \end{columns}
  \vspace{0.6em}
  \begin{itemize}
    \item Bagging laid the foundation for more sophisticated ensemble methods, most notably
          \alert{random forests}.
  \end{itemize}
\end{frame}

%% ════════════════════════════════════════════════════════════════════════════
\section{Random Forests}

\begin{frame}{Random Forests: Beyond Bagging}
  \framesubtitle{Adding feature randomization for decorrelation}
  \begin{itemize}
    \item \alert{Random forests} extend bagging by introducing an additional source of randomness:
          at each node, only a \alert{random subset of features} is considered for splitting.
    \item This feature subsampling decorrelates the trees, reducing the overlap in their
          decision paths.
    \item Because bagging produces trees from overlapping bootstrap samples, those trees tend
          to use similar features at the top levels, leading to \alert{correlation}.
    \item Feature randomization breaks this correlation, leading to greater diversity and
          further variance reduction.
  \end{itemize}
  \begin{block}{Key innovation}
    While bagging randomizes the \alert{data}, random forests randomize both the \alert{data}
    and the \alert{features}, yielding substantially more diverse ensembles.
  \end{block}
\end{frame}

\begin{frame}{Random Forest Algorithm}
  \begin{enumerate}
    \item For $b = 1, \ldots, B$:
      \begin{enumerate}
        \item Draw a bootstrap sample $\mathcal{D}^{(b)}$ of size $n$ from the training set.
        \item Grow a decision tree $T^{(b)}$ on $\mathcal{D}^{(b)}$, with the following
              modification:
              \begin{itemize}
                \item At each node, randomly select $m$ features from the $p$ available features.
                \item Determine the best split using \alert{only those $m$ features}.
                \item Typical choices: $m = \sqrt{p}$ (classification) or $m = p/3$ (regression).
              \end{itemize}
        \item Grow the tree to full depth without pruning (or to a pre-specified minimum
              node size).
      \end{enumerate}
    \item Aggregate predictions: average for regression, majority vote for classification.
  \end{enumerate}
\end{frame}

\begin{frame}{Feature Subsampling Parameter $m$}
  \begin{itemize}
    \item The parameter $m$ controls the strength of feature randomization.
    \item \alert{$m = p$}: no feature randomization, reduces to standard bagging.
    \item \alert{$m = 1$}: maximum randomization, each split considers only one feature.
    \item Intermediate values balance diversity and individual tree quality.
  \end{itemize}
  \begin{block}{Default recommendations}
    \begin{itemize}
      \item \textbf{Classification}: $m = \sqrt{p}$ (e.g., $m = 3$ for $p = 10$).
      \item \textbf{Regression}: $m = p/3$ (e.g., $m \approx 3$ for $p = 10$).
    \end{itemize}
  \end{block}
  \begin{itemize}
    \item Smaller $m$ increases diversity but may reduce individual tree quality.
    \item Larger $m$ reduces diversity but produces stronger individual trees.
    \item In practice, the defaults work well, but $m$ can be tuned via cross-validation.
  \end{itemize}
\end{frame}

\begin{frame}{Why Feature Randomization Works}
  \begin{itemize}
    \item In many datasets, a \alert{few features are much stronger predictors} than the rest.
    \item In bagging, most trees will select these strong features at the top levels,
          making the trees highly similar.
    \item Feature randomization forces the algorithm to consider \alert{alternative features},
          even if they are individually weaker.
  \end{itemize}
  \begin{block}{Decorrelation effect}
    By preventing all trees from relying on the same strong features, random forests create
    more \alert{diverse decision boundaries}, which reduces the correlation among predictions
    and leads to greater variance reduction upon averaging.
  \end{block}
  \begin{itemize}
    \item This is especially beneficial in problems with many correlated features or a few
          dominant predictors.
  \end{itemize}
\end{frame}

\begin{frame}{Deep Trees Without Pruning}
  \begin{itemize}
    \item Unlike classical decision trees, random forests typically grow trees to
          \alert{full depth} without pruning.
    \item Individual trees are allowed to overfit, achieving very low bias but high variance.
    \item The ensemble aggregation then \alert{reduces the variance} without increasing bias.
  \end{itemize}
  \begin{block}{Trade-off}
    \begin{itemize}
      \item Each individual tree is complex and potentially overfit.
      \item The forest as a whole is \alert{regularized by averaging}, not by constraining
            individual trees.
    \end{itemize}
  \end{block}
  \begin{itemize}
    \item This strategy contrasts with boosting, where individual trees are typically kept
          shallow to serve as weak learners.
  \end{itemize}
\end{frame}

\begin{frame}{Out-of-Bag Error in Random Forests}
  \begin{itemize}
    \item As in bagging, each tree in a random forest is trained on a bootstrap sample,
          leaving approximately $36.8\%$ of examples out-of-bag.
    \item OOB predictions for each example are computed using only trees for which that
          example was OOB.
    \item The \alert{OOB error} aggregates these predictions and provides an unbiased estimate
          of generalization performance.
  \end{itemize}
  \begin{block}{Practical advantage}
    OOB error in random forests is nearly as accurate as $k$-fold cross-validation but
    requires \alert{only a single training run}, making it extremely convenient for model
    evaluation and hyperparameter tuning.
  \end{block}
  \begin{itemize}
    \item OOB error typically stabilizes after a few hundred trees and can guide the choice of $B$.
  \end{itemize}
\end{frame}

\begin{frame}{Feature Importance in Random Forests}
  \framesubtitle{Measuring variable relevance}
  \begin{itemize}
    \item Random forests provide a natural mechanism for assessing \alert{feature importance}.
    \item Two common measures:
  \end{itemize}
  \begin{block}{Mean Decrease in Impurity (MDI)}
    For each feature, sum the total reduction in impurity (e.g., Gini or entropy) achieved
    by splits on that feature, averaged over all trees.
    \begin{itemize}
      \item Advantages: fast to compute, available by default.
      \item Disadvantage: biased toward high-cardinality features.
    \end{itemize}
  \end{block}
  \begin{block}{Permutation Importance (Mean Decrease in Accuracy)}
    For each feature, randomly permute its values in the OOB samples and measure the
    decrease in prediction accuracy.
    \begin{itemize}
      \item More reliable and unbiased.
      \item Requires recomputing predictions, more expensive.
    \end{itemize}
  \end{block}
\end{frame}

\begin{frame}{Hyperparameters in Random Forests}
  \renewcommand{\arraystretch}{1.4}
  \begin{center}
  \begin{tabular}{ll}
    \toprule
    \textbf{Hyperparameter} & \textbf{Typical values / guidance} \\
    \midrule
    $B$ (number of trees) & $100$--$1000$; does not cause overfitting \\
    $m$ (features per split) & $\sqrt{p}$ (classif.), $p/3$ (regr.) \\
    Min samples per leaf & $1$ (classif.), $5$ (regr.) \\
    Max depth & Unlimited (full trees) \\
    Bootstrap sample size & $n$ (same as training set) \\
    \bottomrule
  \end{tabular}
  \end{center}
  \vspace{0.4em}
  \begin{itemize}
    \item Random forests are remarkably \alert{robust to hyperparameter choices}.
    \item Default settings often work well out-of-the-box.
    \item The most important parameter is $m$; tuning it can provide modest improvements.
  \end{itemize}
\end{frame}

\begin{frame}{Random Forests: Strengths and Limitations}
  \begin{columns}[t]
    \column{0.48\textwidth}
      \begin{block}{Strengths}
        \begin{itemize}
          \item Excellent \alert{out-of-the-box performance}.
          \item Handles high-dimensional data well.
          \item Provides feature importance estimates.
          \item Robust to overfitting.
          \item Naturally handles mixed feature types.
          \item Easily parallelizable.
        \end{itemize}
      \end{block}
    \column{0.48\textwidth}
      \begin{block}{Limitations}
        \begin{itemize}
          \item Loss of interpretability.
          \item Can be slow for very large datasets.
          \item Less effective than boosting for highly structured problems.
          \item Memory-intensive for large $B$.
        \end{itemize}
      \end{block}
  \end{columns}
  \vspace{0.6em}
  \begin{itemize}
    \item Random forests are among the most widely used ensemble methods in practice,
          valued for their simplicity, robustness, and strong empirical performance.
  \end{itemize}
\end{frame}

%% ════════════════════════════════════════════════════════════════════════════
\section{Boosting}

\begin{frame}{The Boosting Paradigm}
  \framesubtitle{Sequential learning with adaptive emphasis}
  \begin{itemize}
    \item \alert{Boosting} constructs an ensemble \emph{sequentially}, where each new model
          is trained to correct the errors of the current ensemble.
    \item Unlike bagging, which trains models independently in parallel, boosting is inherently
          \alert{adaptive}: later models focus on examples that earlier models found difficult.
    \item The key mechanism: assign \alert{higher weights} to misclassified examples, forcing
          subsequent learners to prioritize them.
  \end{itemize}
  \begin{block}{Core principle}
    By repeatedly focusing attention on hard cases, boosting transforms a sequence of
    \alert{weak learners} (models barely better than random guessing) into a single
    \alert{strong learner} with high accuracy.
  \end{block}
\end{frame}

\begin{frame}{Weak Learners}
  \begin{itemize}
    \item A \alert{weak learner} is a model that performs slightly better than random guessing
          on any distribution over training examples.
    \item For binary classification: a weak learner achieves error rate $\epsilon < 1/2$.
    \item Typical weak learners in boosting:
      \begin{itemize}
        \item \alert{Decision stumps}: trees with only one split (depth 1).
        \item \alert{Shallow trees}: depth 2--6, capturing simple interactions.
      \end{itemize}
    \item The weakness of individual learners is \alert{deliberate}: boosting's strength comes
          from combining many simple models, not from using complex ones.
  \end{itemize}
  \begin{alertblock}{Contrast with bagging}
    Bagging uses \alert{strong, high-variance} base learners (deep trees), while boosting
    uses \alert{weak, low-variance} base learners (shallow trees).
  \end{alertblock}
\end{frame}

\begin{frame}{Boosting vs.\ Bagging: Conceptual Comparison}
  \renewcommand{\arraystretch}{1.5}
  \begin{center}
  \begin{tabular}{lll}
    \toprule
    \textbf{Aspect} & \textbf{Bagging} & \textbf{Boosting} \\
    \midrule
    Training mode & Parallel & Sequential \\
    Base learners & Strong (deep trees) & Weak (shallow trees) \\
    Error focus & Variance reduction & Bias reduction \\
    Weighting & Uniform & Adaptive \\
    Overfitting risk & Low & Moderate to high \\
    Interpretability & Low & Very low \\
    Parallelization & Easy & Difficult \\
    \bottomrule
  \end{tabular}
  \end{center}
\end{frame}

%% ════════════════════════════════════════════════════════════════════════════
\section{AdaBoost}

\begin{frame}{AdaBoost: Adaptive Boosting}
  \framesubtitle{The canonical boosting algorithm}
  \begin{itemize}
    \item \alert{AdaBoost} (Adaptive Boosting) is the most influential and historically
          important boosting algorithm.
    \item Designed for \alert{binary classification} with labels $t_i \in \{-1, +1\}$.
    \item Maintains a set of nonnegative weights $\mathbf{w}^{(j)} = (w_1^{(j)}, \ldots, w_n^{(j)})$
          over training examples, normalized to define a probability distribution.
    \item Initially, all examples are equally important: $w_i^{(0)} = 1/n$.
  \end{itemize}
  \begin{block}{Adaptive reweighting}
    At each stage, AdaBoost increases the weights of \alert{misclassified examples} and
    decreases the weights of correctly classified ones, forcing subsequent learners to
    focus on hard cases.
  \end{block}
\end{frame}

\begin{frame}{AdaBoost Algorithm: Overview}
  \begin{enumerate}
    \item Initialize weights: $w_i^{(0)} = 1/n$ for $i = 1, \ldots, n$.
    \item For $j = 1, \ldots, m$ stages:
      \begin{enumerate}
        \item Train a weak learner $y_j(\mathbf{x})$ by minimizing the \alert{weighted}
              misclassification error under $\mathbf{w}^{(j)}$.
        \item Compute the weighted error:
              \[
              \epsilon^{(j)} = \frac{\sum_{\mathbf{x}_i \in \mathcal{E}^{(j)}} w_i^{(j)}}{\sum_i w_i^{(j)}},
              \]
              where $\mathcal{E}^{(j)}$ is the set of misclassified examples.
        \item Compute the learner weight:
              \[
              \alpha_j = \frac{1}{2}\ln\left(\frac{1 - \epsilon^{(j)}}{\epsilon^{(j)}}\right).
              \]
        \item Update example weights:
              \[
              w_i^{(j+1)} = w_i^{(j)} \cdot \exp\left(-\alpha_j\, t_i\, y_j(\mathbf{x}_i)\right).
              \]
      \end{enumerate}
    \item Final classifier: $H(\mathbf{x}) = \text{sign}\left(\sum_{j=1}^{m} \alpha_j y_j(\mathbf{x})\right)$.
  \end{enumerate}
\end{frame}

\begin{frame}{Weighted Error and Learner Weight}
  \begin{itemize}
    \item The \alert{weighted error} $\epsilon^{(j)}$ measures performance under the current
          distribution $\mathbf{w}^{(j)}$:
          \[
          \epsilon^{(j)} = \frac{\sum_{\mathbf{x}_i \in \mathcal{E}^{(j)}} w_i^{(j)}}{\sum_i w_i^{(j)}}.
          \]
    \item If $\epsilon^{(j)} > 1/2$, the learner performs worse than random guessing.
          In this case, we can flip its predictions to achieve $\epsilon < 1/2$.
    \item The \alert{learner weight} $\alpha_j$ quantifies the contribution of weak learner $y_j$:
          \[
          \alpha_j = \frac{1}{2}\ln\left(\frac{1 - \epsilon^{(j)}}{\epsilon^{(j)}}\right).
          \]
  \end{itemize}
  \begin{block}{Interpretation}
    \begin{itemize}
      \item If $\epsilon^{(j)} \to 0$ (perfect learner), then $\alpha_j \to \infty$.
      \item If $\epsilon^{(j)} = 1/2$ (random), then $\alpha_j = 0$.
      \item Better learners receive higher weights in the final combination.
    \end{itemize}
  \end{block}
\end{frame}

\begin{frame}{Weight Update Rule}
  \begin{itemize}
    \item After training weak learner $y_j$, example weights are updated:
          \[
          w_i^{(j+1)} = w_i^{(j)} \cdot \exp\left(-\alpha_j\, t_i\, y_j(\mathbf{x}_i)\right).
          \]
    \item The product $t_i\, y_j(\mathbf{x}_i)$ is:
      \begin{itemize}
        \item $+1$ if the prediction is correct.
        \item $-1$ if the prediction is incorrect.
      \end{itemize}
    \item Therefore:
      \begin{itemize}
        \item \alert{Correct predictions}: $w_i^{(j+1)} = w_i^{(j)} \cdot e^{-\alpha_j}$ (weight decreases).
        \item \alert{Incorrect predictions}: $w_i^{(j+1)} = w_i^{(j)} \cdot e^{\alpha_j}$ (weight increases).
      \end{itemize}
  \end{itemize}
  \begin{block}{Adaptive emphasis}
    Misclassified examples receive exponentially increasing attention, while correctly
    classified examples gradually fade into the background.
  \end{block}
\end{frame}

\begin{frame}{AdaBoost Weight Distribution Visualization}
  \begin{center}
    \includegraphics[width=.4\linewidth]{figs/Figure14-1.pdf}
  \end{center}
  \begin{itemize}
    \item The figure illustrates how weights evolve across boosting iterations.
    \item Points consistently misclassified accumulate larger weights (shown as larger markers).
    \item These hard cases dominate the objective of later learners, forcing the ensemble to
          refine decision boundaries precisely where classification is most difficult.
  \end{itemize}
\end{frame}

\begin{frame}{Final Ensemble Prediction}
  \begin{itemize}
    \item After $m$ iterations, AdaBoost produces $m$ weak learners $y_1, \ldots, y_m$ with
          corresponding weights $\alpha_1, \ldots, \alpha_m$.
    \item The \alert{final classifier} combines all weak learners through a weighted vote:
          \[
          H(\mathbf{x}) = \text{sign}\left(\sum_{j=1}^{m} \alpha_j\, y_j(\mathbf{x})\right).
          \]
    \item Each weak learner votes with strength proportional to its accuracy (via $\alpha_j$).
    \item Better learners have more influence; weaker learners contribute less.
  \end{itemize}
  \begin{block}{Interpretation}
    The ensemble score $f(\mathbf{x}) = \sum_{j=1}^{m} \alpha_j\, y_j(\mathbf{x})$ measures
    the \alert{confidence} of the prediction. Larger magnitude $|f(\mathbf{x})|$ indicates
    higher confidence.
  \end{block}
\end{frame}

\begin{frame}{AdaBoost Learning Dynamics: Sequence of Weak Learners}
  \begin{center}
    \includegraphics[width=.3\linewidth]{figures/adaboost1.png}
  \end{center}
  \begin{itemize}
    \item The first weak learner produces a coarse separator that inevitably misclassifies
          certain regions.
    \item These misclassified points receive higher weights, shifting the focus for the next learner.
  \end{itemize}
\end{frame}

\begin{frame}{AdaBoost Learning Dynamics: Second Weak Learner}
  \begin{center}
    \includegraphics[width=.3\linewidth]{figures/adaboost2.png}
  \end{center}
  \begin{itemize}
    \item The second weak learner is trained on the reweighted distribution, focusing on
          points that the first learner misclassified.
    \item It produces a different decision boundary that corrects some of the previous errors.
  \end{itemize}
\end{frame}

\begin{frame}{AdaBoost Learning Dynamics: Third Weak Learner}
  \begin{center}
    \includegraphics[width=.3\linewidth]{figures/adaboost3.png}
  \end{center}
  \begin{itemize}
    \item The third weak learner continues this process, refining the decision boundary in
          regions that remain problematic.
    \item Each successive learner addresses the shortcomings of the current ensemble.
  \end{itemize}
\end{frame}

\begin{frame}{AdaBoost Learning Dynamics: Final Decision Boundary}
  \begin{center}
    \includegraphics[width=.5\linewidth]{figures/adaboost4.png}
  \end{center}
  \begin{itemize}
    \item The final ensemble combines all weak learners through weighted voting.
    \item The resulting decision boundary is highly nonlinear and captures complex patterns,
          even though each individual weak learner is simple (e.g., a decision stump or shallow tree).
  \end{itemize}
\end{frame}

\begin{frame}{Why AdaBoost Works: Margin Maximization}
  \begin{itemize}
    \item For a training example $(\mathbf{x}_i, t_i)$, the \alert{margin} is:
          \[
          \text{margin}_i = t_i \cdot f(\mathbf{x}_i) = t_i \sum_{j=1}^{m} \alpha_j\, y_j(\mathbf{x}_i).
          \]
    \item A positive margin indicates correct classification; larger margin $\Rightarrow$
          more confident prediction.
    \item Theoretical analyses show that AdaBoost tends to \alert{maximize the minimum margin}
          of training examples.
  \end{itemize}
  \begin{block}{Generalization via large margins}
    Even after training error reaches zero, AdaBoost continues to increase margins, pushing
    correctly classified points farther from the decision boundary. This improves robustness
    and generalization.
  \end{block}
\end{frame}

\begin{frame}{AdaBoost and Exponential Loss}
  \begin{itemize}
    \item AdaBoost can be interpreted as minimizing the \alert{exponential loss}:
          \[
          L_{\text{exp}}(f) = \sum_{i=1}^{n} \exp(-t_i\, f(\mathbf{x}_i)).
          \]
    \item This loss is convex and differentiable, making it amenable to gradient-based
          optimization.
    \item The weight update rule in AdaBoost corresponds to a coordinate descent step
          minimizing this loss.
  \end{itemize}
  \begin{alertblock}{Sensitivity to outliers}
    The exponential loss is very sensitive to outliers and mislabeled examples: misclassified
    points receive exponentially increasing penalties, which can lead to overfitting on noisy data.
  \end{alertblock}
  \begin{itemize}
    \item This motivates more robust variants like LogitBoost (logistic loss) and gradient boosting
          with other loss functions.
  \end{itemize}
\end{frame}

\begin{frame}{AdaBoost: Regularization and Overfitting}
  \begin{itemize}
    \item Unlike bagging, boosting can \alert{overfit} if the number of iterations $m$ is too large.
    \item As iterations increase, AdaBoost focuses ever more aggressively on hard examples,
          which may include noise or outliers.
    \item Common regularization strategies:
  \end{itemize}
  \begin{block}{Early stopping}
    Monitor performance on a validation set and stop when validation error begins to increase.
  \end{block}
  \begin{block}{Shrinkage (learning rate)}
    Scale each $\alpha_j$ by a factor $\nu \in (0, 1]$:
    \[
    f_j(\mathbf{x}) = f_{j-1}(\mathbf{x}) + \nu \cdot \alpha_j\, y_j(\mathbf{x}).
    \]
    Smaller $\nu$ requires more iterations but produces smoother, more regularized models.
  \end{block}
\end{frame}

%% ════════════════════════════════════════════════════════════════════════════
\section{Gradient Boosting}

\begin{frame}{From AdaBoost to Gradient Boosting}
  \framesubtitle{Generalizing boosting to arbitrary loss functions}
  \begin{itemize}
    \item AdaBoost is elegant but limited to classification with exponential loss.
    \item \alert{Gradient boosting} generalizes the boosting principle to \emph{any differentiable
          loss function}.
    \item Core idea: fit each new weak learner to the \alert{negative gradient} of the loss
          with respect to the current model's predictions.
    \item This is \alert{functional gradient descent}: we move in function space to minimize
          the empirical loss.
  \end{itemize}
  \begin{block}{Universality}
    Gradient boosting works for classification, regression, ranking, and any other task where
    a differentiable loss can be defined.
  \end{block}
\end{frame}

\begin{frame}{Gradient Boosting Algorithm}
  \begin{enumerate}
    \item Initialize the model with a constant prediction:
          \[
          f_0(\mathbf{x}) = \arg\min_{\gamma} \sum_{i=1}^{n} L(t_i, \gamma).
          \]
    \item For $m = 1, \ldots, M$ iterations:
      \begin{enumerate}
        \item Compute the \alert{pseudo-residuals} (negative gradients):
              \[
              r_{im} = -\left.\frac{\partial L(t_i, f(\mathbf{x}_i))}{\partial f(\mathbf{x}_i)}\right|_{f = f_{m-1}}.
              \]
        \item Fit a weak learner $h_m(\mathbf{x})$ to the pseudo-residuals:
              \[
              h_m = \arg\min_h \sum_{i=1}^{n} (r_{im} - h(\mathbf{x}_i))^2.
              \]
        \item Update the model:
              \[
              f_m(\mathbf{x}) = f_{m-1}(\mathbf{x}) + \nu \cdot \alpha_m\, h_m(\mathbf{x}),
              \]
              where $\nu \in (0, 1]$ is the learning rate.
      \end{enumerate}
    \item Final model: $f_M(\mathbf{x})$.
  \end{enumerate}
\end{frame}

\begin{frame}{Pseudo-Residuals: The Gradient Direction}
  \begin{itemize}
    \item At iteration $m$, we have a current model $f_{m-1}(\mathbf{x})$.
    \item To improve the model, we want to move in the direction that \alert{most reduces the loss}.
    \item This direction is the \alert{negative gradient} of the loss:
          \[
          r_{im} = -\frac{\partial L(t_i, f(\mathbf{x}_i))}{\partial f(\mathbf{x}_i)}\bigg|_{f = f_{m-1}}.
          \]
    \item These pseudo-residuals $r_{im}$ represent the \alert{correction} needed for each
          training example.
  \end{itemize}
  \begin{block}{Fitting to gradients}
    By training $h_m$ to predict $r_{im}$, we are effectively fitting a model that approximates
    the \alert{steepest descent direction} in function space.
  \end{block}
\end{frame}

\begin{frame}{Loss Functions in Gradient Boosting}
  \renewcommand{\arraystretch}{1.4}
  \begin{center}
  \begin{tabular}{lll}
    \toprule
    \textbf{Task} & \textbf{Loss function} & \textbf{Gradient} \\
    \midrule
    Regression & Squared: $L(t, f) = (t - f)^2$ & $r = t - f$ \\
    Regression & Absolute: $L(t, f) = |t - f|$ & $r = \text{sign}(t - f)$ \\
    Regression & Huber (robust) & $r = $ clipped residual \\
    Classification & Logistic (cross-entropy) & $r = t - p(t=1 \mid \mathbf{x})$ \\
    Classification & Exponential (AdaBoost) & $r = t \cdot e^{-tf}$ \\
    Ranking & Pairwise loss & $r = $ pairwise gradient \\
    \bottomrule
  \end{tabular}
  \end{center}
  \vspace{0.4em}
  \begin{itemize}
    \item Different loss functions yield different pseudo-residuals, but the gradient boosting
          framework remains the same.
  \end{itemize}
\end{frame}

\begin{frame}{Learning Rate (Shrinkage)}
  \begin{itemize}
    \item The \alert{learning rate} $\nu \in (0, 1]$ controls the step size of each update:
          \[
          f_m(\mathbf{x}) = f_{m-1}(\mathbf{x}) + \nu \cdot \alpha_m\, h_m(\mathbf{x}).
          \]
    \item Smaller $\nu$ produces smoother, more regularized trajectories in function space.
    \item Larger $\nu$ allows faster convergence but increases the risk of overshooting and
          overfitting.
  \end{itemize}
  \begin{block}{Practical guidance}
    \begin{itemize}
      \item Typical values: $\nu \in [0.01, 0.1]$.
      \item Smaller $\nu$ requires more iterations $M$ but often yields better generalization.
      \item Use early stopping on a validation set to select $M$ for a given $\nu$.
    \end{itemize}
  \end{block}
  \begin{itemize}
    \item The combination of small learning rate and early stopping is a \alert{highly effective
          regularization strategy}.
  \end{itemize}
\end{frame}

\begin{frame}{Tree Depth and Interaction Order}
  \begin{itemize}
    \item In gradient boosting, \alert{decision trees} are the standard choice of weak learner.
    \item The depth of these trees controls the \alert{order of interactions} the model can capture:
      \begin{itemize}
        \item \alert{Depth 1 (stumps)}: main effects only, no interactions.
        \item \alert{Depth 2}: pairwise interactions.
        \item \alert{Depth $d$}: up to $d$-way interactions.
      \end{itemize}
    \item Typical recommendation: depth between \alert{3 and 6}.
  \end{itemize}
  \begin{block}{Trade-off}
    \begin{itemize}
      \item Shallow trees: lower variance, require more iterations to fit complex patterns.
      \item Deep trees: higher variance, can overfit, reduce diversity of corrections.
    \end{itemize}
  \end{block}
  \begin{itemize}
    \item Unlike random forests, boosting uses \alert{shallow trees} to maintain the role of
          each tree as a weak learner.
  \end{itemize}
\end{frame}

\begin{frame}{Stochastic Gradient Boosting}
  \framesubtitle{Adding randomness for robustness}
  \begin{itemize}
    \item \alert{Stochastic gradient boosting} introduces randomness at each iteration by
          training each tree on a \alert{random subsample} of the training data.
    \item Subsampling can be applied to:
      \begin{itemize}
        \item \alert{Rows (examples)}: train each tree on a fraction $\eta \in (0, 1]$ of examples.
        \item \alert{Columns (features)}: select a random subset of features for each tree.
      \end{itemize}
    \item This is analogous to the randomization in random forests, but applied within the
          boosting framework.
  \end{itemize}
  \begin{block}{Benefits}
    \begin{itemize}
      \item Reduces correlation between successive trees.
      \item Improves generalization and robustness.
      \item Speeds up training (smaller datasets per tree).
    \end{itemize}
  \end{block}
  \begin{itemize}
    \item Typical values: $\eta \in [0.5, 0.8]$ for row subsampling.
  \end{itemize}
\end{frame}

\begin{frame}{Regularization in Gradient Boosting}
  \begin{itemize}
    \item Gradient boosting is powerful but prone to overfitting without proper regularization.
    \item Key regularization mechanisms:
  \end{itemize}
  \begin{block}{Early stopping}
    Monitor validation error and stop training when it stops improving.
  \end{block}
  \begin{block}{Learning rate (shrinkage)}
    Use small $\nu \in [0.01, 0.1]$ to produce smoother updates.
  \end{block}
  \begin{block}{Tree depth}
    Limit depth to 3--6 to prevent individual trees from overfitting.
  \end{block}
  \begin{block}{Subsampling}
    Use row and/or column subsampling to decorrelate trees.
  \end{block}
  \begin{itemize}
    \item Combining these strategies produces robust, high-performing models.
  \end{itemize}
\end{frame}

\begin{frame}{Why Does Gradient Boosting Work?}
  \framesubtitle{Multiple complementary perspectives}
  \begin{itemize}
    \item \alert{Optimization perspective}: gradient boosting is functional gradient descent,
          minimizing a differentiable loss through iterative additive corrections.
    \item \alert{Bias--variance perspective}: boosting primarily reduces \alert{bias} by
          sequentially correcting systematic errors of simple base learners.
    \item \alert{Statistical perspective}: boosting increases the \alert{classification margin},
          pushing correctly classified examples farther from the decision boundary.
  \end{itemize}
  \begin{block}{Synergy of weak learners}
    Each shallow tree provides a small but meaningful correction; the ensemble builds up
    representational power that no individual tree could achieve. Averaging many corrections
    also provides some variance reduction.
  \end{block}
\end{frame}

%% ════════════════════════════════════════════════════════════════════════════
\section{Modern Gradient Boosting Implementations}

\begin{frame}{XGBoost: Extreme Gradient Boosting}
  \framesubtitle{Engineering optimizations and explicit regularization}
  \begin{itemize}
    \item \alert{XGBoost} (Extreme Gradient Boosting) is one of the most widely used gradient
          boosting implementations.
    \item Key innovations:
  \end{itemize}
  \begin{block}{Second-order optimization}
    XGBoost uses both first-order (gradient) and \alert{second-order (Hessian)} information
    to find better splits:
    \[
    \text{Gain} \propto \frac{G^2}{H + \lambda},
    \]
    where $G$ is the sum of gradients, $H$ is the sum of Hessians, and $\lambda$ is a
    regularization parameter.
  \end{block}
  \begin{block}{Explicit regularization}
    The objective includes regularization terms on the number of leaves and leaf weights:
    \[
    \mathcal{L} = \sum_{i} L(t_i, \hat{t}_i) + \sum_{k} \left[\gamma T + \frac{1}{2}\lambda \|\mathbf{w}\|^2\right].
    \]
  \end{block}
\end{frame}

\begin{frame}{XGBoost: Engineering Optimizations}
  \begin{itemize}
    \item \alert{Column block structure}: features are stored in compressed column format
          (CSC) for efficient split finding.
    \item \alert{Cache-aware access}: data structures are designed to fit in CPU cache,
          reducing memory latency.
    \item \alert{Out-of-core computation}: handles datasets larger than memory by streaming
          data from disk.
    \item \alert{Distributed training}: scales to very large datasets across multiple machines.
    \item \alert{Sparsity-aware splitting}: handles missing values and sparse features efficiently
          by learning default directions for each split.
  \end{itemize}
  \begin{block}{Result}
    XGBoost is both faster and more scalable than traditional gradient boosting implementations
    while maintaining or improving predictive accuracy.
  \end{block}
\end{frame}

\begin{frame}{LightGBM: Light Gradient Boosting Machine}
  \framesubtitle{Histogram-based splitting and leaf-wise growth}
  \begin{itemize}
    \item \alert{LightGBM}, developed by Microsoft, introduces two major innovations:
  \end{itemize}
  \begin{block}{Histogram-based split finding}
    Instead of sorting continuous features at each node, LightGBM bins feature values into
    discrete histograms (typically 255 bins). Split candidates are evaluated on histogram bins,
    reducing computational cost from $O(n\log n)$ to $O(n)$ per feature.
  \end{block}
  \begin{block}{Leaf-wise (best-first) growth}
    Unlike the standard level-wise (breadth-first) strategy, LightGBM always splits the leaf
    that offers the \alert{largest gain}, regardless of depth. This produces asymmetric but
    deeper trees that reduce training loss more rapidly per split.
  \end{block}
  \begin{itemize}
    \item Requires careful tuning of maximum depth to prevent overfitting on small datasets.
  \end{itemize}
\end{frame}

\begin{frame}{LightGBM: Additional Features}
  \begin{itemize}
    \item \alert{Gradient-based One-Side Sampling (GOSS)}: retains all examples with large
          gradients (hard examples) and randomly samples from those with small gradients,
          reducing dataset size while preserving information.
    \item \alert{Exclusive Feature Bundling (EFB)}: bundles mutually exclusive sparse features
          (e.g., one-hot encoded categoricals) into single features, reducing effective
          dimensionality.
    \item Together, these optimizations make LightGBM extremely fast and memory-efficient,
          especially on high-dimensional sparse data.
  \end{itemize}
  \begin{block}{Performance}
    LightGBM is often \alert{10--20× faster} than XGBoost on large datasets while achieving
    comparable or better accuracy.
  \end{block}
\end{frame}

\begin{frame}{CatBoost: Categorical Boosting}
  \framesubtitle{Ordered boosting and native categorical support}
  \begin{itemize}
    \item \alert{CatBoost}, developed by Yandex, addresses two problems:
  \end{itemize}
  \begin{block}{Ordered boosting (target leakage prevention)}
    In standard gradient boosting, gradients are computed using the same examples that train
    the tree, leading to \alert{prediction shift}. CatBoost uses ordered boosting: for each
    example, the gradient is computed from a model trained on a random subset \emph{excluding}
    that example, yielding unbiased gradient estimates.
  \end{block}
  \begin{block}{Native categorical feature handling}
    CatBoost handles categorical features internally using ordered target statistics: for each
    category, the target statistic is computed from a random prefix of the data (respecting a
    permutation), avoiding target leakage while providing informative representations.
  \end{block}
\end{frame}

\begin{frame}{CatBoost: Symmetric Trees}
  \begin{itemize}
    \item CatBoost uses \alert{symmetric (oblivious) trees} as weak learners.
    \item In a symmetric tree, all nodes at the same depth use the \alert{same feature and
          the same split threshold}.
    \item This is more constrained than asymmetric trees (used by XGBoost and LightGBM) but
          has advantages:
  \end{itemize}
  \begin{block}{Benefits of symmetric trees}
    \begin{itemize}
      \item Highly regular computation patterns.
      \item Very fast inference (can be implemented as simple lookup tables).
      \item Less prone to overfitting.
    \end{itemize}
  \end{block}
  \begin{itemize}
    \item CatBoost is particularly fast for \alert{prediction after training}, making it
          attractive for production deployment.
  \end{itemize}
\end{frame}

\begin{frame}{Comparison: XGBoost, LightGBM, CatBoost}
  \renewcommand{\arraystretch}{1.4}
  \begin{center}
  \begin{tabular}{llll}
    \toprule
    \textbf{Feature} & \textbf{XGBoost} & \textbf{LightGBM} & \textbf{CatBoost} \\
    \midrule
    Split finding & Exact / approx & Histogram & Histogram \\
    Tree growth & Level-wise & Leaf-wise & Level-wise \\
    Regularization & Explicit & Explicit & Ordered boosting \\
    Categorical features & Manual encoding & Manual encoding & Native \\
    Speed (training) & Fast & Very fast & Moderate \\
    Speed (inference) & Fast & Fast & Very fast \\
    GPU support & Yes & Yes & Yes \\
    \bottomrule
  \end{tabular}
  \end{center}
  \vspace{0.4em}
  \begin{itemize}
    \item All three are excellent choices; selection depends on dataset characteristics and
          computational constraints.
  \end{itemize}
\end{frame}

%% ════════════════════════════════════════════════════════════════════════════
\section{Practical Considerations}

\begin{frame}{Hyperparameter Tuning in Gradient Boosting}
  \renewcommand{\arraystretch}{1.4}
  \begin{center}
  \begin{tabular}{ll}
    \toprule
    \textbf{Hyperparameter} & \textbf{Typical values / guidance} \\
    \midrule
    Number of trees $M$ & 100--1000; use early stopping \\
    Learning rate $\nu$ & 0.01--0.1 (smaller is better) \\
    Tree depth & 3--6 (shallow trees) \\
    Min samples per leaf & 20--50 (prevents overfitting) \\
    Row subsampling & 0.5--0.8 \\
    Column subsampling & 0.5--1.0 \\
    Regularization $\lambda$ & 0--10 (L2 penalty on weights) \\
    \bottomrule
  \end{tabular}
  \end{center}
  \vspace{0.4em}
  \begin{itemize}
    \item Gradient boosting is more sensitive to hyperparameters than random forests.
    \item Use grid search or Bayesian optimization with cross-validation.
  \end{itemize}
\end{frame}

\begin{frame}{Early Stopping}
  \begin{itemize}
    \item \alert{Early stopping} monitors performance on a validation set during training.
    \item Training stops when validation error stops improving for a specified number of
          iterations (patience).
    \item This is the \alert{most reliable way} to select the number of trees $M$.
  \end{itemize}
  \begin{block}{Procedure}
    \begin{enumerate}
      \item Split data into training and validation sets.
      \item Train the model, evaluating validation error after each iteration.
      \item Stop when validation error has not improved for $p$ consecutive iterations.
      \item Return the model at the iteration with the best validation error.
    \end{enumerate}
  \end{block}
  \begin{itemize}
    \item Typical patience: $p = 10$--$50$ iterations.
    \item Early stopping prevents overfitting and saves computational time.
  \end{itemize}
\end{frame}

\begin{frame}{Ensemble Methods: When to Use Which?}
  \renewcommand{\arraystretch}{1.4}
  \begin{center}
  \begin{tabular}{lll}
    \toprule
    \textbf{Method} & \textbf{Best for} & \textbf{Key strength} \\
    \midrule
    Bagging & High-variance learners & Variance reduction \\
    Random Forests & General-purpose, robust & Out-of-box performance \\
    AdaBoost & Binary classification & Simplicity, interpretability \\
    Gradient Boosting & Structured data, competitions & Bias reduction, accuracy \\
    XGBoost & Large datasets & Speed, scalability \\
    LightGBM & High-dimensional sparse data & Extreme speed \\
    CatBoost & Categorical features & Native categorical support \\
    \bottomrule
  \end{tabular}
  \end{center}
\end{frame}

\begin{frame}{Strengths and Limitations of Ensemble Methods}
  \begin{columns}[t]
    \column{0.48\textwidth}
      \begin{block}{Strengths}
        \begin{itemize}
          \item \alert{State-of-the-art accuracy} on structured data.
          \item Handle complex interactions.
          \item Robust to overfitting (with proper regularization).
          \item Feature importance estimates.
          \item Flexible (work with various losses).
        \end{itemize}
      \end{block}
    \column{0.48\textwidth}
      \begin{block}{Limitations}
        \begin{itemize}
          \item Loss of interpretability.
          \item Computationally expensive (boosting).
          \item Sensitive to hyperparameters (boosting).
          \item Require careful tuning.
          \item Not well-suited for unstructured data (images, text).
        \end{itemize}
      \end{block}
  \end{columns}
  \vspace{0.6em}
  \begin{itemize}
    \item For structured/tabular data, ensembles are the \alert{method of choice}.
    \item For images, text, and audio, deep learning typically outperforms ensembles.
  \end{itemize}
\end{frame}

%% ════════════════════════════════════════════════════════════════════════════
\section{Summary}

\begin{frame}{Summary: Core Concepts}
  \renewcommand{\arraystretch}{1.5}
  \begin{center}
  \begin{tabular}{ll}
    \toprule
    \textbf{Concept} & \textbf{Definition / Key idea} \\
    \midrule
    Ensemble & Combination of multiple models \\
    Bootstrap & Resample with replacement, size $n$ \\
    OOB & Out-of-bag examples ($\approx 36.8\%$) \\
    Bagging & Bootstrap + aggregate (variance $\downarrow$) \\
    Random Forest & Bagging + feature randomization \\
    Boosting & Sequential, adaptive weighting (bias $\downarrow$) \\
    Weak learner & Error $< 1/2$ (shallow trees) \\
    AdaBoost & Exponential loss, weight updates \\
    Gradient boosting & Fit to pseudo-residuals $-\partial L/\partial f$ \\
    Learning rate $\nu$ & Shrinkage, regularization \\
    Early stopping & Stop when validation error stops improving \\
    \bottomrule
  \end{tabular}
  \end{center}
\end{frame}

\begin{frame}{Summary: Key Formulas}
  \begin{columns}[t]
    \column{0.48\textwidth}
      \begin{block}{Bootstrap \& Bagging}
        \begin{itemize}
          \item OOB probability: $(1 - 1/n)^n \to e^{-1}$
          \item Bagging (regr.): $\hat{f} = \frac{1}{B}\sum_b \hat{f}^{(b)}$
        \end{itemize}
      \end{block}
      \vspace{0.4em}
      \begin{block}{AdaBoost}
        \begin{itemize}
          \item Weight: $w_i^{(j+1)} = w_i^{(j)} e^{-\alpha_j t_i y_j(\mathbf{x}_i)}$
          \item $\alpha_j = \frac{1}{2}\ln\frac{1-\epsilon^{(j)}}{\epsilon^{(j)}}$
          \item Final: $H(\mathbf{x}) = \text{sign}(\sum_j \alpha_j y_j)$
        \end{itemize}
      \end{block}
    \column{0.48\textwidth}
      \begin{block}{Gradient Boosting}
        \begin{itemize}
          \item Pseudo-residual: $r_i = -\frac{\partial L}{\partial f(\mathbf{x}_i)}$
          \item Update: $f_m = f_{m-1} + \nu \cdot h_m$
          \item Margin: $\text{margin}_i = t_i f(\mathbf{x}_i)$
        \end{itemize}
      \end{block}
      \vspace{0.4em}
      \begin{block}{Random Forest}
        \begin{itemize}
          \item Features per split: $m \approx \sqrt{p}$ (classif.)
          \item Features per split: $m \approx p/3$ (regr.)
        \end{itemize}
      \end{block}
  \end{columns}
\end{frame}

\begin{frame}{Summary: Methods at a Glance}
  \begin{columns}[t]
    \column{0.48\textwidth}
      \begin{block}{Bagging / Random Forests}
        \begin{itemize}
          \item Parallel training.
          \item Deep trees (low bias, high variance).
          \item Variance reduction through averaging.
          \item OOB error estimation.
          \item Robust, easy to use.
        \end{itemize}
      \end{block}
      \vspace{0.4em}
      \begin{block}{Hyperparameters}
        \begin{itemize}
          \item Random Forest: $B$, $m$, min leaf size.
          \item Boosting: $M$, $\nu$, depth, subsampling.
        \end{itemize}
      \end{block}
    \column{0.48\textwidth}
      \begin{block}{Boosting (AdaBoost, Gradient)}
        \begin{itemize}
          \item Sequential training.
          \item Shallow trees (weak learners).
          \item Bias reduction through corrections.
          \item Adaptive weighting of examples.
          \item Powerful but needs regularization.
        \end{itemize}
      \end{block}
      \vspace{0.4em}
      \begin{block}{Modern Implementations}
        \begin{itemize}
          \item XGBoost: second-order, explicit reg.
          \item LightGBM: histograms, leaf-wise.
          \item CatBoost: ordered boosting, categoricals.
        \end{itemize}
      \end{block}
  \end{columns}
\end{frame}

\begin{frame}{Final Remarks}
  \begin{itemize}
    \item \alert{Ensemble methods} are among the most powerful and widely used techniques in
          machine learning, particularly for structured/tabular data.
    \item They achieve state-of-the-art performance in many real-world applications and
          competitions (e.g., Kaggle).
    \item \alert{Random forests} provide excellent out-of-the-box performance with minimal tuning.
    \item \alert{Gradient boosting} (XGBoost, LightGBM, CatBoost) often achieves the highest
          accuracy but requires careful hyperparameter tuning and regularization.
    \item The key to success: understanding the bias--variance trade-off, selecting appropriate
          base learners, and applying proper regularization.
  \end{itemize}
  \begin{block}{Practical advice}
    Start with random forests for quick exploration; move to gradient boosting for maximum
    performance. Always use cross-validation or OOB error for honest evaluation.
  \end{block}
\end{frame}

%\end{document}
